\subsubsection{Real Exponential and Half-Angle Elements}
The sigma-number-one exponential is evaluated by applying the scalar
exponential separately to the two projector parts defined in
Eq.\ \eqref{eq:functional-calculus}. For the real paravector
\[
\pv{X}=a+\vct{A}\cdot\sigv\in\RPspace,
\]
the vector part is decomposed as
\begin{equation*}
\begin{aligned}
\vct{A}&=\theta \vct{u},\qquad \theta=\lVert \vct{A}\rVert,\\
\vct{u}&=\frac{\vct{A}}{\lVert \vct{A}\rVert}\ \ (\vct{A}\neq 0),\qquad
\vct{u}=\vthree{u_1}{u_2}{u_3},\\
\vct{u}^2&=1.
\end{aligned}
\end{equation*}
When \(\vct{A}=\vct{0}\), the vector exponential is \(1\) and no direction
is needed. The directional formulas below apply to
\(\vct{A}\neq\vct{0}\).
The exponential is therefore given by the projector-based function rule of
Eq.\ \eqref{eq:functional-calculus}:
\begin{align}
\exp(\pm\vct{A}\cdot\sigv)
&= \exp(\pm\theta\,\vct{u}\cdot\sigv)\notag \\
&= \exp(\pm\theta)\pv{\Pi}(+) + \exp(\mp\theta)\pv{\Pi}(-)\notag \\
&= \cosh(\theta) \pm \sinh(\theta)\,\vct{u}\cdot\sigv.
\label{eq:exp-vector}
\end{align}
An equivalent derivation is obtained from the power series by separating
the even and odd powers:
\begin{align*}
\exp(\vct{A}\cdot\sigv)
&= \exp(\theta\,\vct{u}\cdot\sigv)\notag \\
&= \sum_{n\geq0}\frac{\theta^n(\vct{u}\cdot\sigv)^n}{n!}\notag \\
&= \sum_{n\geq0}\frac{\theta^{2n}}{(2n)!}
+ \sum_{n\geq0}\frac{\theta^{2n+1}}{(2n+1)!}\,\vct{u}\cdot\sigv\notag \\
&= \cosh(\theta)+\sinh(\theta)\,\vct{u}\cdot\sigv.
\end{align*}
With
\begin{align*}
\gamma &:= \cosh(\theta),\notag \\
\vct{v} &:= \tanh(\theta)\,\vct{u},
\end{align*}
Eq.\ \eqref{eq:exp-vector} gives the same formula in terms of
\(\cosh(\theta)\), \(\sinh(\theta)\), and their ratio
\(\tanh(\theta)\):
\begin{align*}
\exp(\vct{A}\cdot\sigv) &= \gamma(1+\vct{v}\cdot\sigv),\notag \\
\exp(-\vct{A}\cdot\sigv) &= \gamma(1-\vct{v}\cdot\sigv).
\end{align*}
For the real paravector \(\pv{X}\) from Eq.\ \eqref{eq:real-paravector-x}, the scalar part commutes with the vector part, so
\begin{align*}
\exp(\pv{X}) &= \exp(a)\exp(\vct{A}\cdot\sigv),\notag \\
\exp(\pv{X}^{*}) &= \exp(a)\exp(-\vct{A}\cdot\sigv)=\exp(\pv{X})^{*},\notag \\
\exp(-\pv{X}) &= \exp(-a)\exp(-\vct{A}\cdot\sigv)=\exp(\pv{X})^{-1}.
\end{align*}
The corresponding half-angle element is
\begin{equation}
\pv{L}(\theta,\vct{u})
:=\exp\!\left(\frac{\theta}{2}\,\vct{u}\cdot\sigv\right).
\label{eq:boost-element}
\end{equation}
Its conjugation properties are
\[
\pv{L}^{\Hop}=\pv{L},
\qquad
\conj{\pv{L}}=\pv{L}^{-1}.
\]
The action of this element on real events is derived from the two-sided
event law in Eq.\ \eqref{eq:general-event-congruence}; it is not obtained
from the inverse-pair product in Eq.\ \eqref{eq:general-transform}. The
exponential and its inverse are specified by Eq.\ \eqref{eq:boost-element}. The two
projector weights are exchanged by the conjugate operator:
\begin{align*}
\exp(\vct{A}\cdot\sigv)^{*}
&= \exp(-\theta)\pv{\Pi}(+) + \exp(\theta)\pv{\Pi}(-)\notag \\
&= \gamma(1-\vct{v}\cdot\sigv)\notag \\
&= \exp(-\vct{A}\cdot\sigv)\notag \\
&= \exp\bigl((\vct{A}\cdot\sigv)^{*}\bigr)\notag \\
&= \exp(\vct{A}\cdot\sigv)^{-1}.
\end{align*}

\subsubsection{Real Logarithm of a Paravector}
The logarithm is the inverse problem for the same projector formulas. The
real paravector \(\pv{Y}\) from Eq.\ \eqref{eq:real-paravector-y} is assumed
to satisfy \(b>\lVert\vct{B}\rVert\), so that the two real values in
Eq.\ \eqref{eq:real-paravector-eigenvalues} are positive. The unique ordinary
real logarithm of each positive value is used; this choice is called the
principal logarithm. For \(\vct{B}\neq\vct{0}\), the local direction is
defined by
\begin{equation}
\vct{u}:=\frac{\vct{B}}{\lVert\vct{B}\rVert},
\qquad
\vct{u}^2=1.
\label{eq:log-paravector-axis}
\end{equation}
Thus \(\vct{B}=\lVert\vct{B}\rVert \vct{u}\). The
projector function from Eq.\ \eqref{eq:projectors} is then specialized
immediately after its direction:
\begin{equation}
\pv{\Pi}(\pm):=\tfrac12(1\pm \vct{u}\cdot\sigv).
\label{eq:log-paravector-projectors}
\end{equation}
The eigenvalue construction from
Eq.\ \eqref{eq:real-paravector-eigenvalues} is then specialized to
\(\pv{Y}\). Within this subsection, the two eigenvalues of \(\pv{Y}\) are
denoted locally by \(\lambda(\pm)\):
\begin{align}
\lambda(\pm) &:= b \pm \lVert\vct{B}\rVert,
\label{eq:log-paravector-eigenvalues}\\
\pv{Y} &= \lambda(+)\pv{\Pi}(+) + \lambda(-)\pv{\Pi}(-).\notag
\end{align}
The logarithm is therefore applied separately to the two eigenvalues:
\begin{align}
\log(\pv{Y})
&= \log\!\bigl(\lambda(+)\bigr)\pv{\Pi}(+)\notag\\
&\quad+\log\!\bigl(\lambda(-)\bigr)\pv{\Pi}(-)\notag \\
&= \tfrac12 \log\!\bigl(\lambda(+)\lambda(-)\bigr)
\notag\\
&\quad+ \tfrac12 \log\!\bigl(\lambda(+)/\lambda(-)\bigr)
\,\vct{u}\cdot\sigv.
\label{eq:log-paravector}
\end{align}
When \(\vct{B}=\vct{0}\), the axis is unnecessary and this formula is
reduced directly to \(\log(\pv{Y})=\log(b)\).
For \(\vct{B}\neq0\), if \(\pv{Y}=\exp(\pv{X})\) with \(\pv{X}\) from
Eq.\ \eqref{eq:real-paravector-x} and \(\vct{A}=\theta\vct{u}\), then
\(\lambda(\pm)=\exp(a\pm\theta)\). The relation is expanded as
\begin{align*}
\pv{Y}
&=\exp(\pv{X})\notag \\
&=\exp(a+\vct{A}\cdot\sigv)\notag \\
&=\exp(a)\exp(\vct{A}\cdot\sigv)\notag \\
&=\exp(a)\bigl(\cosh(\theta)+\sinh(\theta)\,\vct{u}\cdot\sigv\bigr)\notag \\
&=b+\lVert\vct{B}\rVert\,\vct{u}\cdot\sigv,
\end{align*}
and the recovered parameters are
\begin{align*}
\theta &= \operatorname{arctanh}\!\bigl(\lVert\vct{B}\rVert/b\bigr), \\
\exp(a) &= \sqrt{b^{2}-\vct{B}^{2}}, \\
a &= \tfrac12 \log\!\bigl(b^{2}-\vct{B}^{2}\bigr), \\
\vct{A} &= \theta\,\vct{u}
= \operatorname{arctanh}\!\bigl(\lVert\vct{B}\rVert/b\bigr)\,
\frac{\vct{B}}{\lVert\vct{B}\rVert}.
\end{align*}
The condition \(b>\lVert\vct{B}\rVert\) gives a positive scalar part and
\(\detf{\pv{Y}}=b^2-\vct{B}^2>0\). In special-relativity language this is
the future-directed timelike region \cite{minkowski1908,einstein1905}.
Thus its principal logarithm is again a real paravector,
\begin{equation*}
\log(\pv{Y})=a+\vct{A}\cdot\sigv.
\end{equation*}

\subsubsection{Circular Exponentials and Half-Angle Elements}
For this subsection, a nonzero real vector \(\vct{A}\) is written with the
local values
\[
\theta:=\|\vct{A}\|,
\qquad
\vct{u}:=\frac{\vct{A}}{\|\vct{A}\|},
\qquad
\vct{A}=\theta\vct{u},
\qquad
\vct{u}^2=1.
\]
Replacing \(\theta\) with \(i\theta\) changes the hyperbolic functions into
the ordinary sine and cosine functions. The resulting sigma-number-two
exponential is therefore called circular:
\begin{align}
\exp(\pm i\vct{A}\cdot\sigv)
&= \exp(\pm i\theta\,\vct{u}\cdot\sigv)\notag \\
&= \exp(\pm i\theta)\pv{\Pi}(+) + \exp(\mp i\theta)\pv{\Pi}(-)\notag \\
&= \cos(\theta) \pm i\sin(\theta)\,\vct{u}\cdot\sigv.
\label{eq:exp-bivector}
\end{align}
Hamilton's quaternion units are obtained by relabeling the same
sigma-number-two basis, without introducing new generators or a new product
\cite{hamilton1844}:
\begin{equation}
\begin{aligned}
\pv{q}_1&:=-\pv{\sigma}_2\pv{\sigma}_3=-i\pv{\sigma}_1,\qquad
\pv{q}_2:=-\pv{\sigma}_3\pv{\sigma}_1=-i\pv{\sigma}_2,\\
\pv{q}_3&:=-\pv{\sigma}_1\pv{\sigma}_2=-i\pv{\sigma}_3,\qquad
\vct{q}:=\vthree{\pv{q}_1}{\pv{q}_2}{\pv{q}_3}=-i\sigv.
\end{aligned}
\label{eq:quaternion-basis}
\end{equation}
The multiplication and intrinsic-operation rules are inherited directly from
Eqs.\ \eqref{eq:sigma-product} and \eqref{eq:bivector-operator-signs}:
\begin{equation}
\begin{aligned}
\pv{q}_n\pv{q}_m
&=-\delta_{nm}+\sum_{k=1}^{3}\epsilon_{nmk}\pv{q}_k,\\
\conj{\pv{q}_n}&=\pv{q}_n,\qquad
\herm{\pv{q}_n}=-\pv{q}_n.
\end{aligned}
\label{eq:quaternion-basis-properties}
\end{equation}
In particular,
\(\pv{q}_1^2=\pv{q}_2^2=\pv{q}_3^2
=\pv{q}_1\pv{q}_2\pv{q}_3=-1\).
The negative-sign exponential may therefore be written as
\begin{equation*}
\exp(-i\vct{A}\cdot\sigv)
=\cos(\theta)+\sin(\theta)\,\vct{u}\cdot(-i\sigv).
\end{equation*}
In the quaternion basis from Eq.\ \eqref{eq:quaternion-basis}, the same
formula is written as
\begin{equation*}
\exp(-i\vct{A}\cdot\sigv)
=\cos(\theta)+\sin(\theta)\,\vct{u}\cdot\vct{q}.
\end{equation*}
For a real angle \(\theta\in\R\) and a unit real vector \(\vct{u}\in\R^3\), the corresponding circular half-angle exponential is
\begin{equation}
\pv{R}(\theta,\vct{u}):=\exp\!\left(\frac{\theta}{2}\,\vct{u}\cdot i\sigv\right).
\label{eq:rotor-element}
\end{equation}
A real scalar coefficient and a purely imaginary vector coefficient are
carried, so \(\pv{R}\in\CPspace\).
Its conjugation properties are
\[
\herm{\pv{R}}\pv{R}=1,
\qquad
\conj{\pv{R}}=\pv{R}.
\]
Equation \eqref{eq:rotor-element} is the circular half-angle element built
from Eq.\ \eqref{eq:exp-bivector}. This element is called a rotor because
its two-sided action produces a spatial rotation. That action is derived in
the Rotation subsection from Eqs.\ \eqref{eq:rotor} and
\eqref{eq:similarity-components}, where the result is reduced to
Rodrigues's rotation formula \cite{doran2003}. The
equivalent quaternion form is obtained from the same sigma-number-two
algebra.

\subsubsection{Exponentials of Complex Paravectors}
The real paravectors are written as
\begin{align*}
\pv{X} &= a + \vct{A} \cdot \sigv \in \RPspace,\notag \\
\pv{Y} &= b + \vct{B} \cdot \sigv \in \RPspace,\notag \\
\pv{Z} &= S + \vct{V} \cdot \sigv \in \CPspace,\notag \\
&= \pv{X} + i \pv{Y},
\end{align*}
where
\begin{equation*}
S=a+ib\in\C,
\qquad
\vct{V}=\vct{A}+i\vct{B}\in\C^3.
\end{equation*}
Because the scalar part \(S\) commutes with \(\vct{V}\cdot\sigv\), the
exponential is separated into scalar and directional factors:
\begin{align*}
\exp(\pv{Z}) &= \exp( S + \vct{V} \cdot \sigv ),\notag \\
&= \exp(S) \exp( \vct{V} \cdot \sigv ).
\end{align*}
The associated star and inverse identities are
\begin{align*}
\exp(\pv{Z}^{*}) &= \exp(S^{*}) \exp( -\vct{V}^{*} \cdot \sigv ),\notag \\
\exp(\pv{Z})^{-1} &= \exp(-S) \exp( -\vct{V} \cdot \sigv ),\notag \\
\exp(\pv{Z}^{*})^{-1} &= \exp(-S^{*}) \exp( \vct{V}^{*} \cdot \sigv ),\notag \\
(\vct{V} \cdot \sigv )^{2} &= \vct{V} \cdot \vct{V} = \vct{V}^{2},\notag \\
&= \vct{A}^{2} - \vct{B}^{2} + 2 i \vct{A} \cdot \vct{B}.
\end{align*}
The key simplification is that all higher powers are reduced to the two basis
elements \(1\) and \(\vct{V}\cdot\sigv\):
\begin{align*}
(\vct{V} \cdot \sigv )^{2n} &= (\vct{V}^{2})^{n},\notag \\
z &= \sqrt{\vct{V}^{2}}\qquad \eqtext{complex scalar},\notag \\
\exp(\vct{V} \cdot \sigv ) &= \sum_{n\geq0} \frac{(\vct{V} \cdot \sigv )^{n}}{n!},\notag \\
&= \sum_{n\geq0} \frac{(\vct{V}^{2})^{n}}{(2n)!}
+ \sum_{n\geq0} \frac{(\vct{V}^{2})^{n}}{(2n+1)!}(\vct{V} \cdot \sigv ),\notag \\
&= \cosh(z) + \sinh(z) \frac{\vct{V} \cdot \sigv}{z}.
\end{align*}
The final expression is independent of the sign chosen for \(z\). When
\(\vct{V}^2=0\), the quotient \(\sinh(z)/z\) is understood by its
continuous value \(1\), and the formula is reduced to
\begin{equation*}
\exp(\pv{Z})=\exp(S)\bigl(1+\vct{V}\cdot\sigv\bigr).
\end{equation*}
